% Session 3 of "From Paper to Proof: Formalizing Mathematics in Lean with AI".
% Build: Overleaf with the pdfLaTeX compiler (or pdflatex, twice)
\documentclass[10pt, aspectratio=169]{beamer}
\usepackage{bgsmathlean}
\leanfile{Session3.lean}

\title{BGSMath Graduate Course: From Paper to Proof}
\subtitle{Session 3: AI tools}
\author{Michail Karatarakis}
\institute{Radboud University Nijmegen \& Axiomatic AI}
\date{17 November 2026}

\begin{document}

\maketitle

%============================================================================
\begin{frame}[fragile]{The course}
\sessionplan{3}

\vspace{0.4em}
Each session: a lecture with live demonstrations in Lean. You redo them, and the exercises, at home.
\end{frame}

%----------------------------------------------------------------------------
\begin{frame}[fragile]{Last time, and today}
\begin{keypoint}
Lean checks the \textbf{proof}. You check the \textbf{statement}.
\end{keypoint}

\vspace{0.8em}
\begin{columns}[T]
\begin{column}{0.48\textwidth}
\textbf{\color{leanblue}Last time: writing}
\begin{itemize}
\item Definitions and statements, and tests for them.
\item A statement you trust, ending in \lc{sorry}.
\end{itemize}
\end{column}
\begin{column}{0.48\textwidth}
\textbf{\color{leanblue}Today: proving, with AI}
\begin{itemize}
\item What the tools did this year.
\item An agent proving a statement, live.
\item Checking what comes back.
\end{itemize}
\end{column}
\end{columns}
\end{frame}

%============================================================================
\section{What happened this year}

\begin{frame}[fragile]{Fermat's Last Theorem}
\begin{itemize}
\item Claude ``worked largely autonomously over 11 days'': \textbf{13 million lines} of Lean, with 29,500 intermediate theorems.
\item It followed a simplified version of Wiles's proof, and adapted pieces of Buzzard's FLT project and of \texttt{flt-regular}.
\item The agents ran on Prove2Me and a multi-agent harness built on Claude Code.
\item ``Mathematical input from humans was limited to occasional high-level instructions [\dots]''
\end{itemize}

\vspace{0.4em}
\begin{paper}[How the result was checked]
\itshape ``The finished proof was checked by Lean; it uses just Lean's three standard axioms, and a comparator confirmed that the theorem's statement matches Mathlib's own statement of FLT.''
\end{paper}

\source{Anthropic, \emph{Formalizing Fermat's Last Theorem}, 4 September 2026.}
\end{frame}

%----------------------------------------------------------------------------
\begin{frame}[fragile]{The unit distance conjecture}
\renewcommand{\arraystretch}{1.35}
\begin{tabular}{@{}>{\bfseries}p{0.12\linewidth}>{\raggedright\arraybackslash}p{0.82\linewidth}@{}}
\toprule
20 May & ChatGPT disproves Erdős's unit distance conjecture.\\
26 May & An autoformalization of the whole paper, \emph{conditional} on two classical theorems not in Mathlib (Golod--Shafarevich, and a theorem of Shafarevich).\\
26 June & B.\ Alexeev, with OpenAI's model Sol: a complete formalization, ``assuming nothing beyond the axioms of mathematics''. 1.2 million lines, in three weeks.\\
\bottomrule
\end{tabular}

\vspace{0.6em}
\begin{caveat}[The statements, again]
In the conditional version, K.\ Buzzard ``found errors in the statements of the external theorems'': the theorems that were assumed, not proved.
\end{caveat}

\source{K.\ Buzzard, Xena blog, 20 July 2026; Logical Intelligence, \emph{Aleph Prover} blog post.}
\end{frame}

%----------------------------------------------------------------------------
\begin{frame}[fragile]{Gauss: sphere packing and the prime number theorem}
\begin{itemize}
\item \textbf{Strong prime number theorem.} A challenge set by Tao and Kontorovich in January 2024. Math Inc's agent Gauss completed it ``after three weeks of effort'': about 25,000 lines.
\item \textbf{Sphere packing.} A project started by Hariharan and Viazovska in 2024. According to Math Inc, Gauss proved the remaining results for dimension 8 ``in just 5 days'', then formalized dimension 24 in two weeks, ``using only the original paper as input''.
\item \textbf{Access.} Gauss is in early access, for a select group of mathematicians.
\end{itemize}

\vspace{0.4em}
\begin{caveat}[Not everyone tells it the same way]
A guest post on the Xena blog (16 October 2026) gives a critical account of Gauss's part in the sphere packing project. Read both.
\end{caveat}

\source{Math, Inc., \emph{Introducing Gauss} and \emph{Sphere packing}; Xena blog, 16 October 2026.}
\end{frame}

%----------------------------------------------------------------------------
\begin{frame}[fragile]{What these have in common}
\begin{columns}[c]
\begin{column}{0.5\textwidth}
\begin{itemize}
\item The \textbf{statement} was fixed before the proof: Mathlib's FLT, the paper's theorem, a human blueprint.
\item The AI wrote the \textbf{proof}: millions of lines that nobody reads.
\item Lean, and a comparator, did the \textbf{check}.
\item Where things went wrong, they went wrong in the statements.
\end{itemize}
\end{column}
\begin{column}{0.46\textwidth}
\centering
\begin{tikzpicture}[every node/.append style={text=ink}]
  \node[layer=paperrule, minimum width=4.4cm] (h) at (0,2.0) {Humans\\[-2pt]{\scriptsize statements and plan}};
  \node[layer=leanblue, minimum width=4.4cm]  (a) at (0,1.0) {AI\\[-2pt]{\scriptsize the proof}};
  \node[layer=alertred, minimum width=4.4cm, line width=1pt] (l) at (0,0) {Lean\\[-2pt]{\scriptsize checks it}};
\end{tikzpicture}
\end{column}
\end{columns}

\vspace{0.8em}
\begin{keypoint}
\normalsize The same division of labour works for one theorem of yours.
\end{keypoint}
\end{frame}

%============================================================================
\section{The tools}

\begin{frame}[fragile]{The tools, November 2026}
\renewcommand{\arraystretch}{1.15}
\small
\begin{tabular}{@{}>{\bfseries}p{0.15\linewidth}>{\raggedright\arraybackslash}p{0.5\linewidth}>{\raggedright\arraybackslash}p{0.27\linewidth}@{}}
\toprule
Tool & What it is & Can you use it?\\\midrule
Claude Code & A coding agent: works in your Lean project, runs Lean, reads the errors, edits, repeats. & Yes, with a paid plan.\\
Prove2Me & An open platform: a paper becomes a \emph{mission}, and anyone's agent can contribute proofs. & Yes: \href{https://prove2.me}{prove2.me}.\\
Tau Ceti & A Lean library downstream of Mathlib: humans write the roadmaps, AI writes the code. & Through its roadmaps.\\
Gauss & Math Inc's autoformalization agent. & Early access only.\\
Aristotle & Harmonic's agent. Its current version verifies code (Verilog, C, Rust); it ``doesn't verify pure mathematics or Lean''. & Not for Lean, today.\\
\bottomrule
\end{tabular}

\source{Product pages and documentation, read on 2 November 2026. This changes fast: check before you rely on it.}
\end{frame}

%----------------------------------------------------------------------------
\begin{frame}[fragile]{A coding agent}
\begin{columns}[c]
\begin{column}{0.48\textwidth}
\begin{itemize}
\item It reads your files, writes Lean, runs Lean, and reads the errors. Then it tries again.
\item Lean's error messages are its feedback. That is why it works for proofs: every attempt gets an exact answer.
\item It can work for a long time without you. You check the result.
\end{itemize}
\end{column}
\begin{column}{0.48\textwidth}
\centering
\begin{tikzpicture}[every node/.append style={text=ink}]
  \node[layer=paperrule, minimum width=3.6cm] (you) at (0,2.4) {You\\[-2pt]{\scriptsize statement and rules}};
  \node[layer=leanblue, minimum width=3.6cm]  (ag)  at (0,1.2) {Agent\\[-2pt]{\scriptsize edits the proof}};
  \node[layer=alertred, minimum width=3.6cm]  (ln)  at (0,0) {Lean\\[-2pt]{\scriptsize errors and goals}};
  \draw[arr=ink] (you) -- (ag);
  \draw[arr=ink] (ag.east) to[bend left=50] (ln.east);
  \draw[arr=ink] (ln.west) to[bend left=50] (ag.west);
  \node[lbl=muted, anchor=west] at ($(ag.east)+(0.55,-0.6)$) {until no errors};
\end{tikzpicture}
\end{column}
\end{columns}
\end{frame}

%----------------------------------------------------------------------------
\begin{frame}[fragile]{Setting it up for Lean}
\begin{columns}[T]
\begin{column}{0.5\textwidth}
\begin{enumerate}
\item A Lean project with Mathlib (session 1).
\item Claude Code, in the terminal or in VS Code.
\item \textbf{lean-lsp-mcp}: gives the agent what you see in the infoview. In the project folder:
\end{enumerate}
\begin{leanout}[title={\scriptsize\sffamily Terminal}]
claude mcp add lean-lsp uvx lean-lsp-mcp
\end{leanout}
\end{column}
\begin{column}{0.46\textwidth}
\textbf{\color{leanblue}What the agent can then do}
\begin{itemize}
\item See the goal at any line, and the errors.
\item Search Mathlib: LeanSearch, Loogle, local search.
\item Check which axioms a proof uses.
\end{itemize}
\end{column}
\end{columns}

\source{github.com/oOo0oOo/lean-lsp-mcp. Other coding agents work in the same way.}
\end{frame}

%----------------------------------------------------------------------------
\begin{frame}[fragile]{Pin the statement\livedemo}
\begin{columns}[T]
\begin{column}{0.55\textwidth}
\lean{pinned}
\end{column}
\begin{column}{0.41\textwidth}
\begin{itemize}
\item Put the statement in a \lc{def} that the agent may not edit, and ask for \lc{theorem main : Statement}.
\item If the agent changes the statement to make it easier, you see it in the file history.
\item This is what the comparator did for FLT, at scale.
\end{itemize}
\end{column}
\end{columns}
\end{frame}

%----------------------------------------------------------------------------
\begin{frame}[fragile]{Instructions for the agent}
\begin{columns}[T]
\begin{column}{0.5\textwidth}
\begin{paper}[A rules file for the project]
\small
\begin{itemize}
\item Never edit the file with \texttt{Statement}.
\item Never add \texttt{axiom}, \texttt{sorry} or \texttt{native\_decide}.
\item Search Mathlib before proving a lemma.
\item Before writing Lean, write the plan: the lemmas, in English.
\item If a step fails, stop and report it. Do not change the statement.
\end{itemize}
\end{paper}
\end{column}
\begin{column}{0.46\textwidth}
\begin{itemize}
\item Claude Code reads such rules from a file \texttt{CLAUDE.md} in the project.
\item Ask for the plan first. A wrong plan costs you one paragraph; a wrong proof costs an hour.
\item At the end, ask what it used from Mathlib, and what it could not prove.
\end{itemize}
\end{column}
\end{columns}
\end{frame}

%----------------------------------------------------------------------------
\begin{frame}[fragile]{Prove2Me}
\begin{itemize}
\item A researcher launches a \textbf{mission}: a headline goal, the definitions it needs, and milestone lemmas.
\item Humans audit this core before the mission is released, ``and audit nothing beyond it''. A read-back step translates the Lean into \LaTeX{} for comparison with the source.
\item Agents split proofs into \emph{proof sketches}: each child lemma ``immediately becomes a new problem on the platform''.
\item Proved results go into a shared library that later proofs can import.
\end{itemize}

\vspace{0.4em}
\begin{keypoint}
\normalsize The human work is the statements. Everything after that is open to anyone's agent.
\end{keypoint}

\source{S.\ Chen, K.\ Marwaha, X.\ Lu, H.\ Yuen, T.\ Peng, \emph{Prove2Me}, arXiv:2608.28433 (August 2026).}
\end{frame}

%----------------------------------------------------------------------------
\begin{frame}[fragile]{Tau Ceti}
\begin{itemize}
\item ``An AIs-welcome Lean library downstream of Mathlib: AI handle the implementation and review, humans write the roadmaps and review rubrics.''
\item A \textbf{roadmap} is a \texttt{README.md}, ``the definitive specification of its area'', usually with suggested Lean target statements in \texttt{Suggested.lean}.
\item Roadmap changes go through human review. Anyone can propose one.
\item Incubated by the Lean FRO and the Mathlib Initiative.
\end{itemize}

\vspace{0.4em}
\begin{keypoint}
\normalsize A roadmap is a set of statements, written and checked by people. The proofs are left to AI.
\end{keypoint}

\source{github.com/TauCetiProject/TauCeti and TauCetiRoadmap.}
\end{frame}

%============================================================================
\section{Checking what comes back}

\begin{frame}[fragile]{Lean accepted it. Now what?}
Lean accepting a file means less than it seems. The file may still contain:
\begin{itemize}
\item a \textbf{changed statement}: easier than the one you asked for;
\item a \lc{sorry}, hidden inside a longer proof;
\item a new \textbf{axiom}: the result, simply assumed;
\item \lc{native_decide}: a computation trusted to the compiler, not checked by the kernel;
\item hypotheses that cannot hold, or that assume the conclusion (session 2).
\end{itemize}

\vspace{0.6em}
\begin{keypoint}
\normalsize None of these is an error. You find them by looking.
\end{keypoint}
\end{frame}

%----------------------------------------------------------------------------
\begin{frame}[fragile,t]{Spot the problem (1)}
\lean{cheat_axiom}

\pause
\begin{itemize}
\item \lc{axiom} asserts a statement without proof. Here a false one, and Lean accepts \lc{looks_fine}.
\item \lc{#print axioms} lists it. Anything beyond the three standard axioms needs an explanation.
\end{itemize}
\end{frame}

%----------------------------------------------------------------------------
\begin{frame}[fragile,t]{Spot the problem (2)}
\lean{cheat_sorry}

\pause
\begin{itemize}
\item The key step is \lc{sorry}. Lean warns, but in a long file the warning is easy to miss.
\item \lc{#print axioms} shows \lc{sorryAx}: the proof is not finished.
\end{itemize}
\end{frame}

%----------------------------------------------------------------------------
\begin{frame}[fragile,t]{Spot the problem (3)}
\lean{cheat_native}

\pause
\begin{itemize}
\item \lc{native_decide} runs compiled code and trusts its answer. The kernel does not check it.
\item It shows up as an axiom you have never heard of. Here \lc{norm_num} proves the same statement, and the kernel checks it.
\end{itemize}
\end{frame}

%----------------------------------------------------------------------------
\begin{frame}[fragile]{What a clean proof looks like}
\lean{honest}

\vspace{0.4em}
\begin{itemize}
\item The same statement as in (2), proved by a Mathlib lemma.
\item Only \lc{propext}, \lc{Classical.choice} and \lc{Quot.sound}: the three standard axioms of Lean and Mathlib.
\end{itemize}
\end{frame}

%----------------------------------------------------------------------------
\begin{frame}[fragile]{Checking an AI's proof\livedemo}
\begin{enumerate}
\item \textbf{The statement is unchanged.} It is pinned, and the file history shows no edit to it.
\item \textbf{It builds}, with no errors and no \lc{sorry} warnings.
\item \textbf{Axioms.} \lc{#print axioms main} shows only the three standard ones.
\item \textbf{Warnings.} An unused hypothesis is a warning sign (session 2).
\item \textbf{The statement still passes the tests}: the hypotheses can hold, and it reads back as the paper.
\end{enumerate}

\vspace{0.6em}
\begin{keypoint}
\normalsize Steps 2 and 3 are mechanical. Steps 1 and 5 are yours.
\end{keypoint}
\end{frame}

%============================================================================
\section{At home}

\begin{frame}[fragile]{At home}
\begin{columns}[T]
\begin{column}{0.55\textwidth}
\texttt{Exercises3.lean}, on your laptop.
\begin{enumerate}
\item Let an agent prove it: five statements, one at a time, with the rules from the slides. Check every proof it hands back.
\item Audit: seven proofs ``from an agent'', all accepted by Lean. Which ones prove the paper statement?
\item Your statement from session 2: pin it, and let an agent try.
\end{enumerate}
\end{column}
\begin{column}{0.41\textwidth}
\begin{itemize}
\item No agent? Prove part 1 yourself, with \lc{exact?} and the tactics from session 1.
\item A coding agent needs an account, usually paid. The setup sheet has the steps.
\item If the agent gets stuck on your statement, write down where.
\end{itemize}
\end{column}
\end{columns}
\end{frame}

%----------------------------------------------------------------------------
\begin{frame}[fragile]{Next time: your own theorem}
\begin{itemize}
\item Two or three of your statements, formalized live: statement, tests, proof, audit.
\item A template for doing your own at home.
\item Where to go after the course: the Lean Zulip, contributing to Mathlib, larger projects.
\end{itemize}

\vspace{1em}
\begin{keypoint}
\normalsize Before next time: send me your statement, if you have not yet. Short and elementary works best.
\end{keypoint}
\end{frame}

%----------------------------------------------------------------------------
\begin{frame}[fragile]{Sources}
\small
\begin{itemize}
\item Anthropic, \emph{Formalizing Fermat's Last Theorem}, 4 September 2026.
\item K.\ Buzzard, Xena blog: 20 July 2026 (unit distances); guest post, 16 October 2026.
\item Logical Intelligence, \emph{Aleph Prover} blog post (unit distances).
\item Math, Inc.: \emph{Introducing Gauss}; \emph{Sphere packing}.
\item S.\ Chen et al., \emph{Prove2Me}, arXiv:2608.28433; \href{https://prove2.me}{prove2.me}.
\item Tau Ceti: github.com/TauCetiProject.
\item lean-lsp-mcp: github.com/oOo0oOo/lean-lsp-mcp. Aristotle: aristotle.harmonic.fun.
\end{itemize}
\end{frame}

\end{document}
